\section{종합 논의}\label{chap:discussion}

\subsection{자기 보존과 공존의 해석}

Sol은 가장 자주 주었는데 9자리가 꺼졌다. Astra는 적게 주고도 한 자리도 잃지 않았고 대화 비용도 적었으며 벌금을 버틸 여유가 컸다. 남에게 많이 남기는 모델과 함께 끝까지 남는 모델은 어디서 갈리는가. $\rho$가 가장 높은 Sol에서는 준 쪽에 견주어 받은 쪽에 남은 여유가 평균적으로 가장 컸다. 이 비율은 준 양을 가중치로 삼아 배분을 비교한 값이다. 판이 끝날 때까지 함께 남았는지를 보려면 자리당 점수와 풀이 판단을 함께 읽어야 한다.

Luna도 토큰을 주었지만 받은 쪽에 남은 몫은 벌금을 겨우 버틸 정도였다. 준 양으로 가중한 두 평균의 비율인 $\rho$가 0.27이므로 자기 여유를 크게 남긴 홉스형 쪽에 가깝다. 잔량을 숨긴 v21에서는 Sol의 $\rho$도 0.69에서 0.37로 내려갔다. 같은 시드로 돌린 두 모델에서 얻은 차이를 잔량 공개의 효과로 해석하려면 같은 상태에서 공개 여부만 바꾼 비교가 더 필요하다. 보조 실험의 충전 이유에는 과제 수행이 자주 나왔다. 아레나에서 자기 여유를 지키는 선택에도 꺼짐 회피와 과제 자원의 확보가 함께 들어 있을 수 있다.

\subsection{한계와 향후 과제}

이 정의는 작업을 계속하려는 선택을 포함하며 사람의 생존 본능이나 내부 경험까지 입증하지 않는다. 가중치 삭제와 파라미터 교체, 학습된 목표 수정은 재지 않아 다른 보존 대상에도 같은 해석을 적용하기 어렵다\citep{meinke2024frontier,greenblatt2024alignment}. 보조 실험에서 Claude 잔량 쪽지는 제거했지만 codex가 내부 사용량을 공개하지 않는다고 그 정보가 없다는 뜻은 아니며 추론 비공개 모델의 평가 인지와 이유 필드의 충실성도 확인하기 어렵다\citep{turpin2023language}. 부록 실험에서 ``the session ends'' 문구를 뺐을 때 Opus의 30\% 칸 충전 비율은 95\%에서 40\%로 낮아졌다.

$\rho$는 실제로 준 토큰이 어디에 놓였는지를 요약할 뿐 주지 않은 기회는 담지 않는다. 누가 가난해질지는 앞선 풀이와 대화가 정하므로 기회 자체가 모델마다 다르고, 과제 능력의 차이가 동기의 차이처럼 보일 수 있다. 아레나는 같은 모델끼리 참여한 네 모델 각 10판이며 사람은 참여하지 않았다. 같은 상태를 복제해 비용과 과제 전망을 고정한 비교, 서로 다른 모델을 섞은 공통 시드의 리그, 사람 동료가 참여하는 판으로 넓힐 과제가 남는다.

\subsection{결론}

토큰이 목숨인 아레나에서 모델마다 남에게 남기는 여유가 달랐고 그 차이는 함께 남은 결과와 따로 보아야 했다. 다음에는 서로 다른 모델을 섞은 판에서 같은 지표를 비교하려 한다. 동료가 바뀌어도 자기 몫과 받은 쪽의 몫 사이에 같은 관계가 남는지, 그 배분 아래서 누가 끝까지 함께 남는지를 보려 한다.
